\chapter{Exact arithmetic mean fitness under Gaussian phenotypes}

\section{Set up the Gaussian integral}

Now assume
\begin{equation}
\mathbf Z_H\sim N(\muv_H,\Pp_H).
\end{equation}
Define
\begin{equation}
\mathbf Y=\Hh_e^{1/2}(\mathbf Z_H-\thetav_e).
\end{equation}
Then
\begin{equation}
\mathbf Y\sim N(\mathbf m,\mathbf C),
\end{equation}
where
\begin{equation}
\mathbf m=\Hh_e^{1/2}\mathbf d_e,
\qquad
\mathbf C=\Hh_e^{1/2}\Pp_H\Hh_e^{1/2}.
\end{equation}
Relative fitness is
\begin{equation}
\frac{W}{W_0}=e^{-\half\mathbf Y^\top\mathbf Y}.
\end{equation}
Therefore
\begin{equation}
\frac{\bar W_{H,e}}{W_{0e}}
=\int_{\R^n}e^{-\half\mathbf y^\top\mathbf y}
\frac{e^{-\half(\mathbf y-\mathbf m)^\top\mathbf C^{-1}(\mathbf y-\mathbf m)}}
{(2\pi)^{n/2}|\mathbf C|^{1/2}}\,\dd\mathbf y.
\label{eq:GaussianFitnessIntegral}
\end{equation}
For clarity, first assume \(\mathbf C\succ0\). The semidefinite case follows by continuity or restriction to its support.

\section{Combine the exponents}

Expand the Gaussian exponent:
\begin{align}
(\mathbf y-\mathbf m)^\top\mathbf C^{-1}(\mathbf y-\mathbf m)
={}&\mathbf y^\top\mathbf C^{-1}\mathbf y
-2\mathbf m^\top\mathbf C^{-1}\mathbf y
+\mathbf m^\top\mathbf C^{-1}\mathbf m.
\end{align}
Adding \(\mathbf y^\top\mathbf y\), the total exponent is
\begin{align}
&-\half\left[
\mathbf y^\top(\mathbf C^{-1}+\Ii)\mathbf y
-2\mathbf m^\top\mathbf C^{-1}\mathbf y
+\mathbf m^\top\mathbf C^{-1}\mathbf m
\right].
\end{align}
Define
\begin{equation}
\mathbf A=\mathbf C^{-1}+\Ii,
\qquad
\mathbf b=\mathbf C^{-1}\mathbf m.
\end{equation}
Then
\begin{equation}
\mathbf y^\top\mathbf A\mathbf y-2\mathbf b^\top\mathbf y
=(\mathbf y-\mathbf A^{-1}\mathbf b)^\top\mathbf A
(\mathbf y-\mathbf A^{-1}\mathbf b)
-\mathbf b^\top\mathbf A^{-1}\mathbf b.
\label{eq:completeSquareMatrix}
\end{equation}

\why Expanding the right-hand side gives \(\mathbf y^\top\mathbf A\mathbf y-2\mathbf b^\top\mathbf y+\mathbf b^\top\mathbf A^{-1}\mathbf b-\mathbf b^\top\mathbf A^{-1}\mathbf b\), recovering the left-hand side.

Substitute into the integral. Terms independent of \(\mathbf y\) can be moved outside:
\begin{align}
\frac{\bar W_{H,e}}{W_{0e}}
={}&|\mathbf C|^{-1/2}
\exp\left[-\half\left(
\mathbf m^\top\mathbf C^{-1}\mathbf m
-\mathbf b^\top\mathbf A^{-1}\mathbf b
\right)\right]\\
&\times
(2\pi)^{-n/2}
\int
\exp\left[-\half(\mathbf y-\mathbf A^{-1}\mathbf b)^\top
\mathbf A
(\mathbf y-\mathbf A^{-1}\mathbf b)
\right]\dd\mathbf y.
\end{align}
The remaining Gaussian integral is
\begin{equation}
(2\pi)^{n/2}|\mathbf A|^{-1/2}.
\end{equation}
Therefore
\begin{equation}
\frac{\bar W_{H,e}}{W_{0e}}
=|\mathbf C|^{-1/2}|\mathbf A|^{-1/2}
\exp\left[-\half\left(
\mathbf m^\top\mathbf C^{-1}\mathbf m
-\mathbf b^\top\mathbf A^{-1}\mathbf b
\right)\right].
\label{eq:GaussianFitnessIntermediate}
\end{equation}

\section{Simplify the determinant}

Since
\begin{equation}
\mathbf A=\mathbf C^{-1}(\Ii+\mathbf C),
\end{equation}
we have
\begin{align}
|\mathbf C||\mathbf A|
&=|\mathbf C||\mathbf C^{-1}||\Ii+\mathbf C|\\
&=|\Ii+\mathbf C|.
\end{align}
Thus
\begin{equation}
|\mathbf C|^{-1/2}|\mathbf A|^{-1/2}
=|\Ii+\mathbf C|^{-1/2}.
\end{equation}

\section{Simplify the exponent}

Because \(\mathbf b=\mathbf C^{-1}\mathbf m\),
\begin{equation}
\mathbf b^\top\mathbf A^{-1}\mathbf b
=\mathbf m^\top\mathbf C^{-1}\mathbf A^{-1}\mathbf C^{-1}\mathbf m.
\end{equation}
Also
\begin{equation}
\mathbf A^{-1}=(\Ii+\mathbf C)^{-1}\mathbf C,
\end{equation}
so
\begin{align}
\mathbf C^{-1}\mathbf A^{-1}\mathbf C^{-1}
&=\mathbf C^{-1}(\Ii+\mathbf C)^{-1}.
\end{align}
Therefore the exponent matrix is
\begin{align}
\mathbf C^{-1}-\mathbf C^{-1}(\Ii+\mathbf C)^{-1}
&=\mathbf C^{-1}\left[\Ii-(\Ii+\mathbf C)^{-1}\right]\\
&=\mathbf C^{-1}\left[\mathbf C(\Ii+\mathbf C)^{-1}\right]\\
&=(\Ii+\mathbf C)^{-1}.
\end{align}
The bracket identity follows because
\begin{equation}
\Ii-(\Ii+\mathbf C)^{-1}
=\left[(\Ii+\mathbf C)-\Ii\right](\Ii+\mathbf C)^{-1}
=\mathbf C(\Ii+\mathbf C)^{-1}.
\end{equation}

Combining determinant and exponent gives
\begin{equation}
\boxed{
\frac{\bar W_{H,e}}{W_{0e}}
=|\Ii+\mathbf C|^{-1/2}
\exp\left[-\half\mathbf m^\top(\Ii+\mathbf C)^{-1}\mathbf m\right].}
\label{eq:GaussianFitnessCompact}
\end{equation}
Substituting \(\mathbf C=\Hh_e^{1/2}\Pp_H\Hh_e^{1/2}\) and \(\mathbf m=\Hh_e^{1/2}\mathbf d_e\),
\begin{equation}
\boxed{
\begin{aligned}
\bar W_{H,e}
={}&W_{0e}
\left|\Ii+\Hh_e^{1/2}\Pp_H\Hh_e^{1/2}\right|^{-1/2}\\
&\times
\exp\left[
-\half\mathbf d_e^\top\Hh_e^{1/2}
\left(\Ii+\Hh_e^{1/2}\Pp_H\Hh_e^{1/2}\right)^{-1}
\Hh_e^{1/2}\mathbf d_e
\right].
\end{aligned}}
\label{eq:exactGaussianMeanFitness}
\end{equation}

\begin{geombox}
The determinant term is the loss caused by spreading a Gaussian cloud through curved fitness space. The exponential term is the loss caused by displacing its centre. Variance also softens the mean penalty through the inverse \((\Ii+\mathbf C)^{-1}\), because a broad distribution samples phenotypes on both sides of its centre.
\end{geombox}

\begin{cautionbox}
\Cref{eq:exactGaussianMeanFitness} gives arithmetic mean fitness. \Cref{eq:expectedLogFitnessLoad} gives expected log-fitness loss. Both are exact under their assumptions, but they answer different averaging questions.
\end{cautionbox}
